%!TEX program = xelatex
% ==============================================================================
% AA CONFIGURAZIONE DI BASE ED ENGINE
% ==============================================================================
% BB CLASSE DEL DOCUMENTO
\documentclass[11pt, a4paper]{article}

% ==============================================================================
% AA LINGUA, FONT E TIPOGRAFIA
% ==============================================================================
% BB GESTIONE FONT XELATEX
\usepackage{fontspec}
\setmainfont{Latin Modern Roman}

% BB LOCALIZZAZIONE LINGUISTICA
\usepackage[italian]{babel}

% BB TIPOGRAFIA E MICRO-ESTENSIONI
\usepackage[protrusion=true]{microtype}
\raggedbottom

% ==============================================================================
% AA LAYOUT DI PAGINA E SPAZIATURA PARAGRAFI
% ==============================================================================
% BB GEOMETRIA DEI MARGINI (A4 verticale con margine sinistro 2.4cm per rilegatura)
\usepackage[left=2.4cm, right=1.6cm, top=2cm, bottom=2cm]{geometry}

% BB CONTROLLO SPAZIATURA TRA PARAGRAFI
\usepackage[skip=3pt plus 1pt minus 1pt]{parskip}

% ==============================================================================
% AA MOTORE DEI COLORI E MACRO LOCALI
% ==============================================================================
% BB CARICAMENTO CORE XCOLOR IN SPAZIO GRAY PURO
\usepackage[gray]{xcolor}
% BB MACRO PER I SIGMA MONOCROMATICI ED ESTENSIONE RIGHE
\newcommand{\sigone}{\symbf{\sigma}_1}
\newcommand{\sigtwo}{\symbf{\sigma}_2}
\newcommand{\sigthree}{\symbf{\sigma}_3}
\newcommand{\sigfour}{\symbf{\sigma}_4}
\newcommand{\sigfive}{\symbf{\sigma}_5}
\newcommand{\sigsix}{\symbf{\sigma}_6}
\newcommand{\sigseven}{\symbf{\sigma}_7}
\renewcommand{\arraystretch}{1.3}

% ==============================================================================
% AA MATEMATICA E SIMBOLI SCIENTIFICI
% ==============================================================================
% BB PACCHETTI MATEMATICI UNICODE
\usepackage{mathtools}
\usepackage[warnings-off={mathtools-overbracket,mathtools-colon}]{unicode-math}
\setmathfont{Latin Modern Math}

% BB METRICHE MATRICI E FORMULE DISPLAY
\setlength{\arraycolsep}{10pt}
\setlength{\abovedisplayskip}{3pt plus 1pt minus 1pt}
\setlength{\belowdisplayskip}{3pt plus 1pt minus 1pt}
\setlength{\abovedisplayshortskip}{0pt plus 1pt}
\setlength{\belowdisplayshortskip}{2pt plus 1pt minus 1pt}
\allowdisplaybreaks

% ==============================================================================
% AA GRAFICA ED ELEMENTI FLOTTANTI
% ==============================================================================
% BB MOTORE DI INCLUSIONE GRAFICA E GESTIONE RIGIDA
\usepackage{graphicx}
\usepackage{float}

% ==============================================================================
% AA TITOLI E STRUTTURA DEL DOCUMENTO (MONOCROMATICI NERI)
% ==============================================================================
% BB PACCHETTO TITLESEC
\usepackage{titlesec}

% BB FORMATTAZIONE TIPOGRAFICA (TITOLI NERI)
\titleformat{\section}{\normalfont\Large\bfseries}{\thesection}{1em}{}
\titleformat{name=\section,numberless}{\normalfont\Large\bfseries}{}{0pt}{}
\titleformat{\subsection}{\normalfont\large\bfseries}{\thesubsection}{1em}{}
\titleformat{name=\subsection,numberless}{\normalfont\large\bfseries}{}{0pt}{}
\titleformat{\subsubsection}{\normalfont\normalsize\bfseries}{\thesubsubsection}{1em}{}
\titleformat{name=\subsubsection,numberless}{\normalfont\normalsize\bfseries}{}{0pt}{}

% BB SPAZIATURA VERTICALE DEI TITOLI
\titlespacing*{\section}{0pt}{1.2ex plus 0.4ex minus 0.2ex}{0.6ex plus 0.2ex}
\titlespacing*{\subsection}{0pt}{1ex plus 0.3ex minus 0.2ex}{0.5ex plus 0.2ex}
\titlespacing*{\subsubsection}{0pt}{0.8ex plus 0.2ex minus 0.1ex}{0.4ex plus 0.1ex}

% ==============================================================================
% AA INTEGRAZIONE AMBIENTE MATLAB
% ==============================================================================
% BB CARICAMENTO PACCHETTO MATLAB
\usepackage{matlab}

% BB ATTIVAZIONE SEGNALIBRI E INDICE MATLAB
\matlabmultipletitles
\matlabhastoc

% BB PERSONALIZZAZIONE BOX CODICE SORGENTE (Margini compatti per pagina orizzontale singola)
\mdfdefinestyle{matlabcode}{%
  topline=false,
  bottomline=false,
  rightline=false,
  leftline=true,
  linewidth=2.5pt,
  linecolor=codeframe,
  backgroundcolor=codebg,
  innertopmargin=3pt,
  innerbottommargin=3pt,
  innerleftmargin=7pt,
  innerrightmargin=6pt,
  skipabove=2pt,
  skipbelow=2pt
}

% ==============================================================================
% AA IPERTESTO, SEGNALIBRI ED ELENCHI
% ==============================================================================
% BB CONFIGURAZIONE HYPERREF E METADATI PDF
\hypersetup{
    hidelinks,
    pdfauthor={Francesco Lo Verde},
    pdftitle={Appendice E3.3.C -- Vettore di posizione 0p6}
}

% BB SEGNALIBRI AVANZATI
\usepackage{bookmark}

% BB FORMATTAZIONE ELENCHI PUNTATI E NUMERATI
\usepackage{enumitem}
\setlist{itemsep=3.5pt plus 1pt minus 1pt, topsep=4pt plus 1pt minus 1pt, partopsep=0pt, parsep=1pt}

\newsavebox{\rotpagebox}

\begin{document}
\thispagestyle{empty}
\begin{lrbox}{\rotpagebox}
\begin{minipage}[t][\textwidth][t]{\textheight}
\vspace{0pt}
\subsection*{Appendice E3.3.C -- Vettore di posizione ${}^{0}p_{6}$}

\begin{matlabcode}
p06 = simplify(p06, 'Steps', 50);
fprintf('Vettore posizione ^(0)p_(6):\n');
disp(p06);
\end{matlabcode}
\begin{matlaboutput}
Vettore posizione ^(0)p_(6):
\end{matlaboutput}

\vspace{10pt}
\begin{matlabsymbolicoutput}
\renewcommand{\arraystretch}{1.85}
\setlength{\arraycolsep}{7pt}
\resizebox{\linewidth}{!}{$\displaystyle
{}^{0}\mathbf{p}_{6}(\mathbf{q}) =
\begin{pmatrix}
a_2 \,\cos \left(q_1 \right)\,\cos \left(q_2 \right)-d_4 \,\sin \left(q_2 +q_3 \right)\,\cos \left(q_1 \right)-d_6 \,\sin \left(q_2 +q_3 \right)\,\cos \left(q_1 \right)\,\cos \left(q_5 \right)+a_3 \,\cos \left(q_1 \right)\,\cos \left(q_2 \right)\,\cos \left(q_3 \right)-a_3 \,\cos \left(q_1 \right)\,\sin \left(q_2 \right)\,\sin \left(q_3 \right)-d_6 \,\sin \left(q_1 \right)\,\sin \left(q_4 \right)\,\sin \left(q_5 \right)+d_6 \,\cos \left(q_1 \right)\,\cos \left(q_4 \right)\,\sin \left(q_2 \right)\,\sin \left(q_3 \right)\,\sin \left(q_5 \right)-d_6 \,\cos \left(q_1 \right)\,\cos \left(q_2 \right)\,\cos \left(q_3 \right)\,\cos \left(q_4 \right)\,\sin \left(q_5 \right)\\
a_2 \,\cos \left(q_2 \right)\,\sin \left(q_1 \right)-d_4 \,\sin \left(q_2 +q_3 \right)\,\sin \left(q_1 \right)-d_6 \,\sin \left(q_2 +q_3 \right)\,\cos \left(q_5 \right)\,\sin \left(q_1 \right)+a_3 \,\cos \left(q_2 \right)\,\cos \left(q_3 \right)\,\sin \left(q_1 \right)+d_6 \,\cos \left(q_1 \right)\,\sin \left(q_4 \right)\,\sin \left(q_5 \right)-a_3 \,\sin \left(q_1 \right)\,\sin \left(q_2 \right)\,\sin \left(q_3 \right)-d_6 \,\cos \left(q_2 \right)\,\cos \left(q_3 \right)\,\cos \left(q_4 \right)\,\sin \left(q_1 \right)\,\sin \left(q_5 \right)+d_6 \,\cos \left(q_4 \right)\,\sin \left(q_1 \right)\,\sin \left(q_2 \right)\,\sin \left(q_3 \right)\,\sin \left(q_5 \right)\\
\frac{d_6 \,\sin \left(q_2 +q_3 \right)\,\sin \left(q_4 +q_5 \right)}{2}-a_3 \,\sin \left(q_2 +q_3 \right)-a_2 \,\sin \left(q_2 \right)-d_4 \,\cos \left(q_2 +q_3 \right)-d_6 \,\cos \left(q_2 +q_3 \right)\,\cos \left(q_5 \right)-\frac{d_6 \,\sin \left(q_4 -q_5 \right)\,\sin \left(q_2 +q_3 \right)}{2}
\end{pmatrix}
$}
\end{matlabsymbolicoutput}

\end{minipage}
\end{lrbox}
\noindent\rotatebox{90}{\usebox{\rotpagebox}}
\end{document}
